\section{分析与讨论}
\label{sec:analysis}
\subsection{时间与空间复杂度}
表\ref{tab:bounds} 讨论从给定源点进行一次完整搜索。
图本身占 $O(n+m)$ 空间，距离数组占 $O(n)$；表中的队列空间单列。
标准库版本使用懒删除，其堆中可能保留 $O(m)$ 个条目，
与每个顶点仅一个在堆条目的减键实现存在实质差别。

\begin{table}[H]\centering\small
\caption{队列策略的代价与适用条件}\label{tab:bounds}
\begin{tabularx}{\linewidth}{lXXl}\toprule
策略 & 时间上界或分析说明 & 适用条件 & 队列空间\\\midrule
标准库懒删除 & $O(n+m\log(m+1))$ & 非负权；可重复条目 & $O(m+1)$\\
索引二叉堆 & $O((n+m)\log(n+1))$ & 非负权；可直接减键 & $O(n)$\\
段树堆 & $O((n+m)\log(n+1))$ & 顶点固定映射为叶子 & $O(n)$\\
Fibonacci 堆 & $O(m+n\log(n+1))$，摊还 & 非负权；直接定位节点 & $O(n)$\\
PBDS 两种堆 & 依库堆的具体摊还界，不在本项目中证明 & GCC 扩展与有效句柄 & $O(n)$\\
Radix heap & 按键位宽 $B$ 计为 $O(m+nB)$ 的桶迁移分析 & 非负整数；单调键与容量限制 & $O(n+B)$\\
FIFO SPFA & 与成功松弛和重复入队有关，存在坏例 & 本实验只使用非负权 & 依实际条目数\\
SLF+LLL & 启发式不能保证 Dijkstra 式效率 & 本实验只使用非负权 & $O(n)$\\\bottomrule
\end{tabularx}\end{table}

对 Fibonacci 堆，所有弧扫描和成功减键总计 $O(m)$，每个顶点至多一次有效取最小值，
合计 $O(n\log n)$。稀疏图上该界与二叉堆同阶，但前者每次操作涉及多条链接维护，
理论差异不保证常数因子更小。
Radix heap 的分析利用桶重分配时编号降低；若最大有效键界为 $U$，可令
$B=\lceil\log_2(U+1)\rceil$。文献中的专门 radix 设计可给出
$O(m+n\log C)$ 的 Dijkstra 界（$C$ 为边权上界）\cite{radix}；
本项目没有证明其有限 32 桶和压缩句柄实现满足文献全部设定，因此不直接照搬更紧的界。

本轮没有采集各算法的峰值内存曲线、硬件缓存计数、成功松弛次数或桶重分配次数。
空间复杂度分析来自数据结构，不能当作已经完成的内存实测。

\subsection{为何 Fibonacci 堆没有最快}
本轮 NY、BAY 的平均出度只有约 2.78 和 2.49，属于稀疏路网。
Fibonacci 堆减少了减键的渐近代价，但仍须执行合并同度根、维护链表和级联剪切。
其节点虽然一次性分配在数组中，操作依然沿节点链接跳转；二叉堆的大部分状态则是紧凑数组。
较多的链接操作及访问模式可能解释部分差距，这是基于实现结构的解释，
尚需性能计数器与操作计数实验验证，不能把缓存未命中当作本轮已经观察到的事实。

稠密族中边扫描工作随 $n^2$ 增长，堆维护占比可能下降，因此某些队列的差距缩小。
本轮最大稠密规模为 $n=4096$，不存在“在足够大规模最终谁胜出”的实验结论。
随机图的最大顶点规模虽较大，但它属于另一种稀疏度，不能拿来代替大规模稠密实验。

\subsection{Radix 与 SPFA 结果的意义}
在非负整数路网上，radix heap 利用键的单调性，减少比较堆中的上浮和下沉工作，
成为本轮七种 Dijkstra 队列中最快的一种。
它仍要求正确的键范围与句柄容量，不能直接替换一般实数键或非单调任务中的优先队列。

SLF+LLL 在两图上比所有堆版本更快，FIFO 版本却慢一个数量级以上。
这说明标签修正算法高度依赖调度顺序，不能只以“都是 SPFA”推断性能。
两种实现还分别使用重复条目与在队标记，故差距也不能只归因于 SLF/LLL 一个因素。
随机图上的良好表现不构成对抗输入下的效率保证。

\subsection{模型拟合与测量限制}
\plotfigure{complexity-normalized}{按预设模型归一化的 CPU 时间。曲线并非严格常数；该图用于观察模型偏差与有限规模效应，不用于证明复杂度。}{fig:normalized}
本轮复杂度图是有限个离散规模上的计时。归一化曲线的变化既可能来自
算法操作数变化，也可能来自缓存、编译器、频率或输入可达性变化。
RMS 小只能说明所选范围内模型贴合，不能排除另一种模型或推出紧确界。

三轮随机/复杂度实验不足以稳定估计尾部分布；路网只有单批次、无逐查询分布，
两个路网规模也不足以作可靠的规模拟合。WSL 共享宿主硬件，本轮未固定 CPU 亲和性
或测量宿主负载；绝对时间需复验。多实现对拍和单元测试仍不等于形式化证明。

\subsection{后续改进与底稿结论}
后续应增加路网独立批次与逐查询记录，扩大经校验的图规模；
记录松弛、过期条目、合并与桶迁移次数，并用性能计数器检验常数开销的解释。
若研究真正的点对点查询，应统一加入目标提前终止，再考虑双向搜索或预处理算法。

本轮结果说明，理论摊还优势与实际运行速度需要分别评价，实测结论也应受输入范围约束。
\FloatBarrier
